\chapter{Related Work, Limitations, Future Directions}
\label{ch:related}

This final chapter places \Deppy{} in its scientific context,
revisits the limitations from \Cref{ch:soundness}, and outlines
work that would extend the system.

\section{Related work}

\subsection{Lean, Coq, Agda --- the proof assistants}

The reference points are the mature dependent-type proof
assistants: \Lean{}~4 \citep{LeanFour}, \Coq{} (recently rebranded
as Rocq), and \Agda{}.  These systems share a \CIC-flavoured
core, support powerful tactic languages, and have rich library
ecosystems.

\Deppy{}'s proof language is small compared to the tactic
ecosystems of these systems.  Mathlib \citep{LeanCommunityMathlib}
provides tens of thousands of lemmas.  \Deppy{} does not aim to
replicate this; its lemma library lives in user-supplied sidecars.
What \Deppy{} adds is the bridge from idiomatic Python source to
the proof world.  We see this less as a competitor to Mathlib than
as a way to reach Python-domain users who would not otherwise write
Mathlib lemmas.

\subsection{Refinement-type systems}

Refinement-type systems (\textsc{Liquid Haskell}, \textsc{F*},
\textsc{Dafny}, \textsc{Whiley}, \textsc{Stainless}) attach
predicates to types and discharge them with SMT.  This is a closer
relative of \Deppy{}'s arithmetic-Z3 path: the
\texttt{RefinementType} dataclass is the obvious refinement form,
and Z3-discharge of side conditions is the standard refinement
mechanic.

The key difference is the lift to homotopy / cubical reasoning.
Refinement types are about \emph{predicates} on values;
\Deppy{} also handles \emph{paths between values} (transport, Kan
fillers, univalence).  This is what lets \Deppy{} say things about
duck-typing equivalence and monkey-patching that refinement-type
systems cannot.

\subsection{Cubical proof assistants}

Cubical \Agda{} \citep{ABCFHL, ABCHFL2017}, \texttt{redtt}, \texttt{cooltt}, and
similar implement cubical type theory directly.  These systems are
the canonical implementation of the cubical metatheory we draw on.

\Deppy{} is \emph{inspired} by cubical type theory but is not an
implementation of it.  In particular:
\begin{itemize}[leftmargin=1.7em]
  \item \Deppy{}'s interval, path types, and Kan fillers are
    metatheoretic abstractions captured in dataclasses, not
    operational primitives that compute.
  \item The cubical Agda computation rules (composition, Glue
    transport) are not implemented; \Deppy{} models them as proof
    obligations checked structurally.
  \item Univalence is admitted, not proved by an interval-based
    \texttt{ua}.
\end{itemize}

We borrow the \emph{vocabulary} (paths, fillers, glue, cocycles)
and the \emph{conceptual framework} (types as spaces, equivalences
as identifications) without claiming the operational equivalence.
This is honest: a user who cares about univalence's computational
content should use Cubical Agda; a user who wants to verify
properties of an existing Python library has a different
problem and \Deppy{} addresses it.

\subsection{Python verification frameworks}

Several frameworks attempt static analysis of Python:
\begin{itemize}[leftmargin=1.7em]
  \item \textsc{mypy} / \textsc{pyright} --- static type checking
    against PEP 484 type annotations \citep{PEP484}.  Type-level
    only; no proof.
  \item \textsc{Crosshair} --- symbolic-execution-based
    counter-example finder.  Closer to our counterexample-search
    pass than to our verifier.
  \item \textsc{Nagini} --- a Viper-based deductive verifier for
    Python.  Requires Python rewritten with explicit annotations,
    not sidecar files; targets a smaller subset.
  \item \textsc{Hypothesis} --- property-based testing, not
    verification.  We engage Hypothesis-style sampling in the
    \texttt{computational} discharge mode of
    \texttt{SidecarVerificationBackend}.
\end{itemize}

\Deppy{}'s positioning: a sidecar-based verifier with a strong
proof language and a dual deppy-kernel + Lean backend.  No other
framework we know of has the dual-target structure with an
explicit cubical metatheory.

\subsection{Specification and contract languages}

\textsc{ACSL} (for C), \textsc{JML} (for Java), \texttt{contracts.py}
(for Python), and PEP 316 (Programming by Contract for Python)
all attach pre/post conditions and invariants to functions.

\Deppy{}'s sidecar files are similar in spirit but more flexible:
specs can carry foundations, axioms, and PSDL proofs; verify
blocks are first-class; the trust accounting is explicit.

\subsection{Property-based testing}

\textsc{QuickCheck} (Haskell) and \textsc{Hypothesis} (Python) test
properties on randomly-generated inputs.  \Deppy{}'s
\texttt{SidecarVerificationBackend} runs computational tests when
the kernel cannot discharge an axiom, with the result trust level
\texttt{LLM\_CHECKED} (computational evidence is weaker than SMT or
kernel verification).

\section{Limitations recap}
\label{sec:limits-recap}

We summarise the limitations from \Cref{ch:soundness}:

\begin{itemize}[leftmargin=1.7em]
  \item Type system: flat \texttt{Int} carriers in the Lean
    preamble; no universe polymorphism; refinement predicates not
    auto-checked.
  \item Body translation: gaps for try/except, while, generator,
    dynamic dispatch, C-extension boundaries.
  \item Z3: transcendentals, nested quantifiers, IEEE 754 not
    handled.
  \item Lean: no Mathlib dependency by default; opaque foundation
    Props for non-Z3 cases; \texttt{def := sorry} placeholders for
    method bodies.
  \item PSDL: no partial-proof credit; no tactic suggestion mode.
  \item Cohomology: advisory only, not part of the trust chain.
  \item Pipeline: single-pass, no parallelism by default, hard-coded
    \Lean{} timeout.
\end{itemize}

\section{Future directions}

We outline the most promising extensions, in rough order of
expected impact.

\subsection{Concrete typed preamble}

The largest soundness improvement would be to drop the flat-Int
preamble and emit \emph{typed} class structures.  For each binder
class \texttt{Point}, the preamble would emit:

\begin{lstlisting}[language=lean4]
structure Point where
  x : Int
  y : Int
\end{lstlisting}

with method bodies \texttt{Point.distance := fun p q => sqrt(\dots)}
that the elaborator can unfold substantively.  This requires:
\begin{itemize}[leftmargin=1.7em]
  \item type inference from \texttt{inspect.signature} that
    distinguishes nominal types from \texttt{Int};
  \item the body translator emitting Lean tied to the typed
    structure;
  \item the mechanization preamble using
    \texttt{Class.field} accessors instead of
    \texttt{dot\_field}.
\end{itemize}

The benefit is that PSDL-emitted tactics like
\texttt{exact Foundation\_Real\_sqrt\_nonneg} elaborate against
real arithmetic claims, not opaque Props.  The cost is more
elaboration ambiguity; some method-resolution patterns will need
explicit type ascription.

\subsection{Mathlib integration}

Optional Mathlib import would let foundations involving real
numbers (\texttt{sqrt}, \texttt{sin}, \texttt{exp}) discharge as
real lemmas rather than admit as axioms.  The escape hatch
\texttt{lean\_imports: | import Mathlib} already lets users opt in;
the next step is a curated set of foundation $\to$ Mathlib lemma
mappings the certificate emits automatically.

\subsection{Refinement-type elaboration}

The \texttt{RefinementType} dataclass currently carries a
\texttt{predicate} string that the kernel does not unfold.  An
extension would parse the predicate, type-check it against the
local context, and elaborate the refinement
$\{x : T \mid \varphi(x)\}$ as Lean's \texttt{Subtype}.

\subsection{Try/except translation}

The body translator could handle \texttt{try / except} by emitting
a \Lean{} \texttt{Except} or \texttt{Option} monad.  The exception
fibre then becomes the \texttt{Except.error} branch; tactics can
case-split on the result.  This requires:
\begin{itemize}[leftmargin=1.7em]
  \item the body translator's return-type inference to detect
    exception escape paths;
  \item the verify-block pipeline to engage the
    \texttt{ConditionalEffectWitness} term automatically;
  \item PSDL's \texttt{try / except} to compile against the
    monadic encoding.
\end{itemize}

\subsection{Cohomology obstructions}

The cohomology layer currently produces trivial cocycles in the
absence of safety predicates.  A meaningful extension would:
\begin{itemize}[leftmargin=1.7em]
  \item detect when the safety pipeline produces non-trivial
    preconditions on a function;
  \item emit substantive C¹ cocycles witnessing precondition
    propagation through call sites;
  \item emit substantive C² coherences;
  \item run \texttt{Bicomplex.compute\_total\_cohomology} to find
    cross-level $H^2$ obstructions;
  \item refuse certification when an $H^2$ obstruction is non-zero.
\end{itemize}

This would turn the cohomology layer from an advisory diagnostic
into a real verification step --- one that detects
inter-procedural coherence failures the per-function checks miss.

\subsection{LLM iteration loop}

\PSDL{} has affordances for LLM iteration (\texttt{show\_goal()},
\texttt{hint("\dots")}, the \texttt{lean("\dots")} escape hatch).
A natural extension is an iterative loop:
\begin{enumerate}[leftmargin=1.7em]
  \item user/LLM proposes a PSDL proof;
  \item compiler runs; on failure, surface the goal state at the
    failure point;
  \item LLM proposes a refined PSDL block;
  \item repeat until \texttt{certify\_or\_die} passes.
\end{enumerate}

The compiler already produces \texttt{show\_goal\_requested}
flags and per-error diagnostics; assembling the loop is the
remaining work.

\subsection{Multi-target}

The \texttt{deppy.lean} package emits \Lean{}~4.  Other targets are
plausible:
\begin{itemize}[leftmargin=1.7em]
  \item \Coq{} / Rocq: the same trust chain, different tactic
    syntax.
  \item \Agda{}: similar.
  \item Z3-only target: emit the entire certificate as an SMT-LIB
    file (where this is possible); skip \Lean{}.
  \item No-target dry run: the kernel verifies but no certificate
    is emitted.  Useful for incremental checking in a CI pipeline.
\end{itemize}

\subsection{Incremental verification}

\texttt{deppy/pipeline/incremental.py} sketches a per-function
incremental verifier keyed on \texttt{source\_hash}.  The
challenges are:
\begin{itemize}[leftmargin=1.7em]
  \item invalidation when a foundation changes (every theorem
    citing it must re-verify);
  \item invalidation when an axiom changes;
  \item handling pipeline outputs that depend on the entire module
    (cohomology, bicomplex).
\end{itemize}

\subsection{IDE integration}

A language server protocol implementation that:
\begin{itemize}[leftmargin=1.7em]
  \item provides hover-over types for PSDL identifiers;
  \item highlights \texttt{Structural} leaks at edit time;
  \item runs the \Lean{} round-trip in the background and surfaces
    elaboration errors at the PSDL line responsible.
\end{itemize}

\section{Closing remarks}

\Deppy{} occupies a particular niche: it is a verification
framework for Python libraries that someone has already written
and is unwilling to retype.  This niche has been under-served by
the proof-assistant community and over-served by the
property-based-testing community; we offer something in between.

The cubical-duck-type metatheory is the project's intellectual
organising principle.  We do not claim that Python is
\emph{secretly} a cubical type theory.  We claim that Python's
characteristic phenomena --- duck typing, \texttt{isinstance},
exceptions, \texttt{\_\_getattr\_\_}, monkey-patching --- have
\emph{natural} cubical readings, and that exploiting these
readings produces a proof language that is both usable and
expressive.

\PSDL{} is the user-facing surface.  It accepts literal Python and
emits both kernel-checkable proof terms and Lean tactics.  When
the kernel and \Lean{} agree, we have a certified proof.  When
they disagree, we have a diagnostic the user can act on.

The trust accounting is the project's ethical commitment.  Every
admission, every \texttt{Structural} leaf, every untranslated
body is recorded and counted.  The \texttt{CERTIFIED} verdict is
strict; the gates are conservative.  A user reading a \Deppy{}
certificate sees both the positive ``\Lean{} kernel accepted''
side and the negative ``these foundations admit'' side, and can
make their own judgment.

The work that remains is significant: the typed preamble, the
Mathlib integration, the cohomology obstructions, the LLM
iteration loop, the incremental pipeline.  We have left these as
future work in the present document but each is articulated
clearly enough that contributing maintainers can pick one up.

What we hope to have demonstrated, with this monograph, is that
the verification of imported Python modules is approachable --- by
LLMs and humans both --- and that a sidecar-based,
cubical-duck-type-flavoured proof system is a coherent way to
attack it.  The next step is to use it.

\section{Acknowledgements}

The implementation that this monograph documents is the
\Deppy{} project's mainline source tree.  The \texttt{AUDIT.md}
file at the top of the repository was the principal design-rationale
input.  Where this document and \texttt{AUDIT.md} disagree, the
source code wins.

The \HoTT{} book \citep{HoTT}, the cubical type theory papers
\citep{CCHM, ABCFHL, ABCHFL2017}, the \Lean{} reference manual
\citep{LeanRefManual}, and the \Lean{}-mathlib community
\citep{LeanCommunityMathlib} provided the proof-theoretic
background.  The \Zthree{} paper \citep{Z3} provided the SMT
foundation.  PEP 484 \citep{PEP484} and PEP 544 \citep{PEP544} are
the Python language references the verifier targets.

This monograph itself is generated alongside the source tree at
\texttt{docs/psdl\_treatise/}.  The build instructions are at the
top of \texttt{main.tex} and in the local Makefile.
